\documentclass[11pt,a4paper]{article}
\usepackage[margin=2.8cm]{geometry}
\usepackage{amsmath,amssymb,amsthm,mathtools,mathrsfs}
\usepackage[T1]{fontenc}
\usepackage{microtype}
\usepackage{booktabs}
\usepackage[colorlinks=true,linkcolor=blue,citecolor=blue,urlcolor=blue,hypertexnames=false]{hyperref}

\newtheorem{theorem}{Theorem}[section]
\newtheorem{proposition}[theorem]{Proposition}
\newtheorem{lemma}[theorem]{Lemma}
\newtheorem{hypothesis}[theorem]{Hypothesis}
\theoremstyle{remark}
\newtheorem{remark}[theorem]{Remark}

\newcommand{\R}{\mathbb R}
\newcommand{\Hh}{\mathcal H_0}
\newcommand{\Om}{\Omega}
\newcommand{\Q}{\mathcal Q}
\newcommand{\E}{\mathbb E}
\newcommand{\CE}{\mathrm{CE}}
\DeclareMathOperator{\supp}{supp}
\DeclareMathOperator{\spec}{spec}
\DeclareMathOperator{\dist}{dist}

\title{A rigorous weak-coupling strategy for the two remaining inputs\\
in the continuum $P(\varphi)_2$ cutoff problem}
\author{}
\date{8 August 2026}

\begin{document}
\maketitle

\begin{abstract}
The corrected continuum LPPL theorem in \texttt{lppl.tex} has exactly two open inputs: a
Hamiltonian gap uniform in the spatial cutoff and a uniform local second moment of the
interaction density.  This note gives a single source-backed strategy for both inputs in a
precisely specified weak-coupling regime.  The required source estimates are the cutoff-uniform
cluster and polynomial-moment estimates of Glimm--Jaffe--Spencer.  We state those estimates in the
form actually needed, prove the Euclidean-slab-to-ground-state transfer, derive an explicit
uniform bound for the local second moment, prove density of the cylinder vectors, and then derive
the physical Hamiltonian gap by the spectral theorem.  Thus the remaining work is not a new
spectral-gap argument: it is a finite source-scope audit and a careful normalization check.
No stochastic-quantisation gap or change of Dirichlet form occurs.
\end{abstract}

\section{Regime and exact target}

Fix a real semibounded polynomial
\[
 \mathscr P(\xi)=\sum_{n=0}^{p}a_n\xi^n,
 \qquad p\ \text{even},\quad a_p>0,
\]
a free mass $m_0>0$, and a coupling $\lambda\ge0$.  Set
$\mathcal G_0:=C_c^\infty(\R;[0,1])$.  For $g\in\mathcal G_0$, put
\[
 K_g=H_0(m_0)+\lambda W(g),\qquad
 W(f)=\int_\R f(x):\mathscr P(\varphi_0(x)):\,dx,
\]
and let $E_g=\inf\spec K_g$, $H_g=K_g-E_g$.  Its normalized positive ground
state is $\Om_g$ and $P_g=|\Om_g\rangle\langle\Om_g|$.

The weak-coupling regime is the following quantified statement:
\begin{equation}\label{eq:wc}
 0\le \frac{\lambda}{m_0^2}<\varepsilon_{\mathscr P},                 \tag{WC}
\end{equation}
where $\varepsilon_{\mathscr P}>0$ is the constant furnished by the norm-decreasing
cluster expansion in \cite[Thms.~1.1.7--1.1.8]{GJS}.  The statement ``weak coupling'' below
always means \eqref{eq:wc}; it never means an unspecified neighbourhood.

Define the exact test-function class
\[
 \mathcal F_2:=\{f\in C_c^\infty(\R;\R):\|f\|_\infty\le1
 \text{ and }\supp f\subset I\text{ for some interval }I\text{ with }|I|\le2\}.
\]
The two estimates to prove are
\begin{align}
 \spec H_g&\subseteq\{0\}\cup[\gamma_*,\infty)
 &&\text{for every }g,                                             \tag{G}\label{eq:G}\\
 \sup_{g\in\mathcal G_0,\ f\in\mathcal F_2}
 \|\lambda W(f)\Om_g\|^2&<\infty .                               \tag{M$_2$}\label{eq:M2}
\end{align}

\section{The published cutoff-uniform input}

Let $\mu_0^{(2)}$ be the massive free Euclidean field measure on $\mathcal S'(\R^2)$.
For every measurable compactly supported $h:\R^2\to[0,1]$, define
\[
 d\nu_h=Z_h^{-1}\exp\left[-\lambda\int_{\R^2}
 h(z):\mathscr P(\Phi(z)):\,dz\right]d\mu_0^{(2)}.
\]
Glimm--Jaffe--Spencer use unit-square localizations and a norm
$\|\cdot\|_{\CE}$ on sharp-time polynomial observables.  For a localized monomial of local
degrees $N(\Delta)$ and smearing kernel $w$, their norm has the form
\begin{equation}\label{eq:CEnorm}
 \|Q\|_{\CE}=K_1
 \left(\prod_\Delta K_2^{2N(\Delta)}N(\Delta)!\right)\|w\|_2,       \tag{2.1}
\end{equation}
and a general polynomial is assigned the sum of the norms of its localized pieces.

The source input is exactly:

\begin{hypothesis}[GJS cutoff-uniform estimates]\label{hyp:CE}
Under \eqref{eq:wc}, there are $C_{\CE}<\infty$ and $m_{\CE}>0$, independent of
$h$, such that:
\begin{enumerate}
\item for every sharp-time polynomial $Q$ in the class of
\cite[(1.1.5)--(1.1.9)]{GJS},
\begin{equation}\label{eq:CEmoment}
 |\nu_h(Q)|\le C_{\CE}\|Q\|_{\CE};                                \tag{CE1}
\end{equation}
\item if $Q_1,Q_2$ have unit-square localizations separated by a strip of width $d$, then
\begin{equation}\label{eq:CEcluster}
 |\nu_h(Q_1Q_2)-\nu_h(Q_1)\nu_h(Q_2)|
 \le C_{\CE}e^{-m_{\CE}d}\|Q_1\|_{\CE}\|Q_2\|_{\CE}.             \tag{CE2}
\end{equation}
\end{enumerate}
\end{hypothesis}

Estimate \eqref{eq:CEcluster} is Theorem~1.1.7 of \cite{GJS}, whose statement explicitly says
that its constants are independent of $h,Q_1,Q_2$.  Estimate \eqref{eq:CEmoment} is
Theorem~1.1.8; in the proof it is the $B=1,n=0$ case of Theorem~1.1.11, whose error estimate is
explicitly uniform in $h$.  The allowed cutoff in \cite[(1.8)]{GJS} is an arbitrary compactly
supported function $0\le h\le1$, not merely a rectangle.  These three source clauses are the
load-bearing scope checks.

\section{From a Euclidean slab to the cutoff ground state}

Let $\Om_0$ be the free Fock vacuum.  For $T>0$ define
\[
 h_{T,g}(t,x)=\mathbf1_{[-T,T]}(t)g(x),\qquad \nu_{T,g}:=\nu_{h_{T,g}}.
\]
The characteristic function is admissible in the measurable-cutoff formulation of
\cite[(1.8)]{GJS}.  In fact, immediately after (1.8) the source itself takes
$h_j(t,x)=g(x)\chi_{[-j,j]}(t)$; no limiting argument from smooth time cutoffs is required.

\begin{lemma}[one-time slab transfer]\label{lem:one}
For every bounded time-zero multiplication operator $A$,
\[
 \nu_{T,g}(A)=\langle\psi_{T,g},A\psi_{T,g}\rangle,
 \qquad
 \psi_{T,g}:=\frac{e^{-TH_g}\Om_0}{\|e^{-TH_g}\Om_0\|},
\]
and $\psi_{T,g}\to\Om_g$ in Hilbert norm as $T\to\infty$.
\end{lemma}

\begin{proof}
The Feynman--Kac--Nelson formula gives
\[
 \nu_{T,g}(A)=
 \frac{\langle\Om_0,e^{-TK_g}Ae^{-TK_g}\Om_0\rangle}
      {\langle\Om_0,e^{-2TK_g}\Om_0\rangle}.
\]
The factors $e^{-TE_g}$ cancel, giving the first identity.  Positivity improvement in
$Q$-space implies
$c_g:=\langle\Om_g,\Om_0\rangle>0$.  The spectral theorem gives
\[
 e^{-TH_g}\Om_0=c_g\Om_g+e^{-TH_g}(1-P_g)\Om_0
 \longrightarrow c_g\Om_g.
\]
Indeed, the squared norm of the second term is
\[
 \int_{(0,\infty)}e^{-2TE}\,d\mu_g^\perp(E),
 \qquad
 \mu_g^\perp(B):=
 \langle(1-P_g)\Om_0,\mathbf1_B(H_g)(1-P_g)\Om_0\rangle,
\]
which tends to zero by dominated convergence.  Thus even isolation of the ground-state
eigenvalue has not been used.  The norm tends to $c_g$, proving the assertion without any
gap assumption.
\end{proof}

\begin{lemma}[two-time slab transfer]\label{lem:two}
Let $A,B$ be bounded time-zero multiplication operators and $t\ge0$.  If the insertions are at
Euclidean times $-t/2$ and $t/2$, then
\[
 \lim_{T\to\infty}\nu_{T,g}(A(-t/2)B(t/2))
 =\langle\Om_g,Ae^{-tH_g}B\Om_g\rangle.
\]
The assertion remains true for polynomial $A,B$ if
\begin{equation}\label{eq:UI}
 \sup_{T\ge t/2+1}\nu_{T,g}(|A(\pm t/2)|^4+|B(\pm t/2)|^4)<\infty. \tag{3.1}
\end{equation}
\end{lemma}

\begin{proof}
For bounded insertions, FKN and cancellation of the scalar energy give the exact quotient
\[
 \frac{\langle u_{T,t},Ae^{-tH_g}Bu_{T,t}\rangle}
      {\|e^{-TH_g}\Om_0\|^2},
 \qquad u_{T,t}=e^{-(T-t/2)H_g}\Om_0.
\]
Both $u_{T,t}$ and $e^{-TH_g}\Om_0$ tend to $c_g\Om_g$, so the quotient has the asserted limit.

For an unbounded polynomial $A$, put $A_R=A\mathbf1_{\{|A|\le R\}}$, and similarly for $B$.
The bounded result applies to $A_R,B_R$.  By H\"older and Markov,
\begin{align*}
 \nu_{T,g}(|AB|\mathbf1_{\{|A|>R\}})
 &\le \nu_{T,g}(|A|^4)^{1/4}\nu_{T,g}(|B|^4)^{1/4}
       \nu_{T,g}(|A|>R)^{1/2}\\
 &\le R^{-2}\nu_{T,g}(|A|^4)^{3/4}\nu_{T,g}(|B|^4)^{1/4}.
\end{align*}
The analogous bound holds with $A,B$ exchanged.  Hypothesis \eqref{eq:UI} makes both errors
uniformly $O(R^{-2})$.  Moreover $A_R\Om_g\to A\Om_g$ and
$B_R\Om_g\to B\Om_g$ in norm, since the corresponding second moments are finite.  The
contraction property of $e^{-tH_g}$ therefore gives convergence of the truncated matrix
elements on the right.  First take $T\to\infty$, then $R\to\infty$.
\end{proof}

By \eqref{eq:CEmoment}, condition \eqref{eq:UI} holds for every fixed GJS polynomial: its fourth
power is again a finite sum of observables of the form (1.1.5), and its $\CE$-norm is finite.

\section{Explicit derivation of the local second moment}

\begin{proposition}[CE1 implies (M$_2$)]\label{prop:M2}
Hypothesis~\ref{hyp:CE}(CE1) implies \eqref{eq:M2}.  One admissible bound is
\begin{equation}\label{eq:Mexplicit}
 \|\lambda W(f)\Om_g\|^2
 \le 32\lambda^2 C_{\CE}K_1K_2^{4p}(2p)!
       \left(\sum_{n=0}^p|a_n|\right)^2.                            \tag{4.1}
\end{equation}
\end{proposition}

\begin{proof}
Write $I$ for an interval of length at most two containing $\supp f$.  It meets at most four
unit spatial cells.  Expanding the square gives
\[
 W(f)^2=\sum_{n,m=0}^pa_na_m
 \iint f(x)f(y):\varphi_0(x)^n::\varphi_0(y)^m:\,dx\,dy.             \tag{4.2}
\]
Split each integral according to the at most $4^2=16$ ordered pairs of cells.  In each localized
piece the kernel is a restriction of $f\otimes f$, so
\[
 \|f\otimes f\|_2=\|f\|_2^2\le |I|\|f\|_\infty^2\le2.
\]
The total local degree is $n+m\le2p$.  Therefore
\[
 \prod_\Delta N(\Delta)!\le(n+m)!\le(2p)!,
 \qquad
 \prod_\Delta K_2^{2N(\Delta)}=K_2^{2(n+m)}\le K_2^{4p}.
\]
Enlarge $K_2$ to $\max\{K_2,1\}$ if necessary.  Using \eqref{eq:CEnorm}, summing the sixteen
pieces and then $n,m$ proves
\begin{equation}\label{eq:observable-norm}
 \|W(f)^2\|_{\CE}
 \le32K_1K_2^{4p}(2p)!\left(\sum_{n=0}^p|a_n|\right)^2.             \tag{4.3}
\end{equation}
Thus CE1 bounds $\nu_{T,g}(\lambda^2W(f)^2)$ by the right side of
\eqref{eq:Mexplicit}, uniformly in $T,g,f$.

Here the equality
\[
 \|\lambda W(f)\psi_{T,g}\|^2=\nu_{T,g}(\lambda^2W(f)^2)
\]
is also justified for the unbounded observable: apply the bounded FKN identity to
\[
 W(f)^2\wedge R:=\min\{W(f)^2,R\},
\]
and let $R\uparrow\infty$ by monotone convergence.  The CE1 bound makes the resulting quantity
finite.

By Lemma~\ref{lem:one}, $\psi_{T,g}\to\Om_g$.  Since $W(f)$ is closed, the quadratic form
$\xi\mapsto\|W(f)\xi\|^2$ is lower semicontinuous in Hilbert norm.  Hence
\[
 \|\lambda W(f)\Om_g\|^2
 \le\liminf_{T\to\infty}\|\lambda W(f)\psi_{T,g}\|^2,
\]
and \eqref{eq:Mexplicit} follows.
\end{proof}

The numerical constant $32$ is deliberately crude; its purpose is to show that no dependence on
the position of $I$, the profile $g$, or the derivatives of $f$ is hidden in the argument.

\section{Explicit derivation of the physical Hamiltonian gap}

Let $\mathscr C$ be the complex unital algebra of finite polynomials in compactly smeared real
sharp-time fields; its self-adjoint part $\mathscr C_{\rm sa}$ consists of the real-valued
polynomial cylinders.

\begin{lemma}[density of polynomial cylinder vectors]\label{lem:dense}
For each fixed $g$,
\[
 \overline{\{F\Om_g:F\in\mathscr C\}}=\Hh.
\]
Consequently $\{(1-P_g)F\Om_g:F\in\mathscr C\}$ has dense span in
$(1-P_g)\Hh$.
\end{lemma}

\begin{proof}
In $Q$-space let $d\mu_g=\Om_g^2d\mu_0$; strict positivity of $\Om_g$ makes
$U:L^2(\mu_g)\to L^2(\mu_0)$, $UF=F\Om_g$, unitary.

The $n$-point bound following Theorem~1.1.8 of \cite{GJS}, followed by the slab transfer and
truncation, gives for every compactly smeared real field $X=\varphi_0(f)$ a constant
$b_f<\infty$ such that
\[
 \E_{\mu_g}[X^{2n}]\le(2n)!b_f^{2n}\qquad(n\ge0).                  \tag{5.1}
\]
Thus, for $0<a<b_f^{-1}$,
\begin{equation}\label{eq:exp-moment}
 \E_{\mu_g}[e^{a|X|}]
 \le2\E_{\mu_g}[\cosh(aX)]
 \le2\sum_{n=0}^\infty (ab_f)^{2n}<\infty.                         \tag{5.2}
\end{equation}

We now prove the finite-dimensional density assertion rather than invoke it.  Let
$Y=(\varphi_0(f_1),\ldots,\varphi_0(f_d))$ and let $\rho$ be its law under $\mu_g$.
By H\"older and \eqref{eq:exp-moment}, there is $\alpha>0$ such that
\begin{equation}\label{eq:vector-exp}
 \int_{\R^d}\exp\!\left(2\alpha\sum_{j=1}^d|y_j|\right)d\rho(y)<\infty. \tag{5.3}
\end{equation}
Suppose $u\in L^2(\rho)$ is orthogonal to every polynomial.  On the connected tube
\[
 \mathcal T_\alpha:=\{z\in\mathbb C^d:|\operatorname{Re}z_j|<\alpha
 \text{ for }1\le j\le d\}
\]
define
\[
 L(z)=\int_{\R^d}u(y)e^{z\cdot y}\,d\rho(y).
\]
Cauchy--Schwarz and \eqref{eq:vector-exp} give an integrable majorant on every compact subset
of $\mathcal T_\alpha$; the same is true after multiplication by any monomial in $y$.
Consequently differentiation under the integral is valid and $L$ is holomorphic on
$\mathcal T_\alpha$.  Orthogonality gives
\(
 \partial^\beta L(0)=\int u(y)y^\beta\,d\rho(y)=0
\)
for every multi-index $\beta$.  Hence the Taylor series of $L$ at zero vanishes.  The identity
theorem (applied successively in the $d$ variables) gives $L=0$ throughout
$\mathcal T_\alpha$.  In particular,
\[
 \int u(y)e^{i\xi\cdot y}\,d\rho(y)=0\qquad(\xi\in\R^d).
\]
Uniqueness of the Fourier transform of the finite complex measure $u\,d\rho$ implies
$u=0$ $\rho$-almost everywhere.  Thus polynomials are dense in $L^2(\rho)$.

Choose a countable dense family $(f_j)_{j\ge1}$ of real Schwartz functions.  The coordinate
maps $\varphi_0(f_j)$ generate the completed Gaussian cylinder sigma-algebra; because
$\Om_g>0$ $\mu_0$-almost everywhere, $\mu_g$ and $\mu_0$ are equivalent and have the same
completed null sets.  If $\mathcal F_N=\sigma(\varphi_0(f_1),\ldots,\varphi_0(f_N))$, the
martingale convergence theorem gives
\(
 \E_{\mu_g}[Z\mid\mathcal F_N]\to Z
\)
in $L^2(\mu_g)$ for every $Z\in L^2(\mu_g)$.  Each conditional expectation is an $L^2$
function of the first $N$ coordinates and hence, by the preceding paragraph, is an $L^2$ limit
of polynomials in those coordinates.  This proves density in $L^2(\mu_g)$.  Applying $U$ proves
the first assertion, and applying $1-P_g$ proves the second.
\end{proof}

\begin{proposition}[CE2 implies the uniform physical gap]\label{prop:gap}
Hypothesis~\ref{hyp:CE}(CE2) implies
\[
 \spec H_g\subseteq\{0\}\cup[m_{\CE},\infty)
 \qquad\text{for every }g.
\]
\end{proposition}

\begin{proof}
Fix $F\in\mathscr C_{\rm sa}$ and decompose it into finitely many localized monomials as
in \cite[(1.1.9)]{GJS}.  Translate one copy by Euclidean time $t$.  For every pair of localized
summands, the two collections of unit squares are separated by a horizontal strip of width at
least $t-2$.  Summing CE2 over the pairs therefore gives, for $t\ge2$ and $T\ge t/2+1$,
\begin{equation}\label{eq:cluster-time}
 |\nu_{T,g}(F(-t/2)F(t/2))-
   \nu_{T,g}(F(-t/2))\nu_{T,g}(F(t/2))|
 \le C_{\CE}e^{2m_{\CE}}\|F\|_{\CE}^2e^{-m_{\CE}t}.                \tag{5.4}
\end{equation}
The fourth-moment version of CE1 verifies the uniform-integrability hypothesis of
Lemma~\ref{lem:two}.  Letting $T\to\infty$ in \eqref{eq:cluster-time} yields
\begin{equation}\label{eq:semigroup-decay}
\begin{aligned}
0&\le
 \langle(1-P_g)F\Om_g,e^{-tH_g}(1-P_g)F\Om_g\rangle\\
 &\le C_Fe^{-m_{\CE}t}\qquad(t\ge0).
\end{aligned}
\tag{5.5}
\end{equation}
For $0\le t<2$, the left side is at most
\(
 \|(1-P_g)F\Om_g\|^2\le C_{\CE}\|F^2\|_{\CE}
\)
by CE1 and the slab transfer.  Thus one may take
\[
 C_F=e^{2m_{\CE}}C_{\CE}
 \max\{\|F\|_{\CE}^2,\|F^2\|_{\CE}\},
\]
which is independent of $g$.  In particular, the decay exponent is independent of both $g$
and $F$.

For completeness, fix $0<\varepsilon<m_{\CE}$ and let
\[
 R_\varepsilon=
 \mathbf1_{[0,m_{\CE}-\varepsilon]}(H_g)(1-P_g).
\]
The spectral theorem and \eqref{eq:semigroup-decay} give
\[
 e^{-(m_{\CE}-\varepsilon)t}\|R_\varepsilon(1-P_g)F\Om_g\|^2
 \le C_Fe^{-m_{\CE}t}.
\]
Multiplication by $e^{(m_{\CE}-\varepsilon)t}$ and $t\to\infty$ show that
$R_\varepsilon(1-P_g)F\Om_g=0$.  Lemma~\ref{lem:dense} and boundedness of
$R_\varepsilon$ give $R_\varepsilon=0$: indeed, every $G\in\mathscr C$ is
$G=(G+G^*)/2+i(G-G^*)/(2i)$, a complex linear combination of two elements of
$\mathscr C_{\rm sa}$.  As $\varepsilon\downarrow0$, these projections increase strongly to
$\mathbf1_{[0,m_{\CE})}(H_g)(1-P_g)$; hence that projection is zero, which proves the claimed
spectral inclusion.
\end{proof}

\section{Execution order and falsification gates}

The work should be carried out in this order.

\begin{enumerate}
\item \textbf{Freeze the theorem.}  State the first unconditional result only for the family
$K_g=H_0(m_0)+\lambda W(g)$ satisfying \eqref{eq:wc}.  Do not write ``away from criticality'';
that is a larger, unproved region.

\item \textbf{Complete the source-scope audit.}  Check in the proof of
\cite[Thms.~1.1.8 and 1.1.11]{GJS} that every occurrence of the cutoff is bounded only by
$0\le h\le1$.  Record the page and equation for each use.  Failure criterion: any constant
depending on the diameter, shape, derivatives, or time extent of $h$ invalidates the transfer.

\item \textbf{Check conventions.}  Match the free covariance, Wick mass, coefficient $\lambda$,
and the energy shift between \cite{GJS} and \texttt{lppl.tex}.  If the source Wick mass is
$\widetilde m\ne m_0$, use the identity
\[
 :\!\xi^n\!:\,_{\widetilde m}
 =\sum_{k=0}^{\lfloor n/2\rfloor}
   \frac{n!}{2^k k!(n-2k)!}
   \left(\frac{1}{2\pi}\log\frac{\widetilde m}{m_0}\right)^k
   :\!\xi^{,n-2k}\!:\,_{m_0}.                                    \tag{6.1}
\]
This follows by comparing the generating functions
$\exp(t\xi-C_m(0)t^2/2)$ and using the finite covariance difference
$C_{m_0}(0)-C_{\widetilde m}(0)=(2\pi)^{-1}\log(\widetilde m/m_0)$.  Recompute all
lower-degree coefficients before applying \eqref{eq:wc}; merely renaming the Wick mass is not
valid.

\item \textbf{Prove the two transfers.}  Insert Lemmas~\ref{lem:one}--\ref{lem:two} into the
cutoff framework.  The positivity overlap $\langle\Om_g,\Om_0\rangle>0$ and the fourth-moment
tail bound must be written explicitly; weak convergence alone is insufficient for polynomial
insertions.

\item \textbf{Discharge (M$_2$) first.}  Verify the observable-norm calculation
\eqref{eq:observable-norm} directly from \cite[(1.1.7)--(1.1.9)]{GJS}.  This is a finite calculation and is
the earliest decisive test of the route.

\item \textbf{Discharge (G).}  Transfer CE2, prove Lemma~\ref{lem:dense}, and apply the spectral
projection argument in Proposition~\ref{prop:gap}.  Do not cite the infinite-volume particle
mass theorem as a substitute: what is needed is the finite-cutoff estimate uniform in $g$.

\item \textbf{Assemble.}  Substitute $\gamma=m_{\CE}$ and the constant from
\eqref{eq:Mexplicit} into Theorem~5.3 of \texttt{lppl.tex}.  The conclusion is full-net local-norm
convergence for the weak-coupling family, not merely convergence of Schwinger distributions.
\end{enumerate}

\section{Routes not to pursue}

\begin{enumerate}
\item An unweighted stochastic-quantisation LSI has the wrong Dirichlet form and the wrong
generator.  It is unnecessary in the weak-coupling route above.
\item The physical gap cannot control the ultraviolet second moment through a double commutator;
the resulting coincident-point Wick square diverges.
\item The infinite-volume mass gap of \cite[Thm.~1.1.2]{GJS} does not by itself imply a gap for
every inhomogeneous spatial cutoff.  The load-bearing input is the cutoff-uniform finite-volume
bound \eqref{eq:CEcluster}.
\item No analytic continuation in $\lambda$ should be used to enlarge \eqref{eq:wc} without a
uniform lower bound on the gap along the continuation path; analyticity of finite-volume
quantities does not prevent a full-net gap from closing.
\end{enumerate}

\section{Precise status}

\begin{theorem}[closure conditional only on the quoted GJS estimates]\label{thm:status}
Assume Hypothesis~\ref{hyp:CE} with the stated cutoff scope.  Then the two hypotheses (G) and
(M$_2$) of \texttt{lppl.tex} hold with constants independent of
$g\in C_c^\infty(\R;[0,1])$.  Consequently the full spatial-cutoff ground-state net converges in
norm on every local algebra in the weak-coupling regime \eqref{eq:wc}, with all conclusions of
Theorem~5.3 of \texttt{lppl.tex}.
\end{theorem}

\begin{proof}
Proposition~\ref{prop:M2} proves (M$_2$), Proposition~\ref{prop:gap} proves (G), and
Theorem~5.3 of \texttt{lppl.tex} performs the assembly.
\end{proof}

The only remaining uncertainty in this strategy is documentary rather than conceptual: the
published $O(1)$ in the moment estimate must be certified as uniform over the entire cutoff class
used in the slab transfer.  Theorem~1.1.11 says ``uniformly in $h$'', and Theorem~1.1.8 is its
$B=1,n=0$ specialization, so the expected outcome is positive.  If the line-by-line audit confirms
that reading, the two formerly open inputs are simultaneously closed in the regime
\eqref{eq:wc}.  No claim is made outside that regime.

\enlargethispage{2\baselineskip}
\begin{thebibliography}{9}
\bibitem{GJS}
J.~Glimm, A.~Jaffe, and T.~Spencer,
\emph{The Wightman axioms and particle structure in the $\mathscr P(\varphi)_2$ quantum field
model}, Ann. of Math. \textbf{100} (1974), 585--632;
\href{https://doi.org/10.2307/1970959}{doi:10.2307/1970959}.
\end{thebibliography}

\end{document}
